\begin{tikzpicture}[
  x=1mm,y=1mm,figure picture,
  state/.style={circle,draw=FigureStroke,fill=white,line width=.75pt,
    minimum size=8.5mm,inner sep=0pt,font=\figurefont\footnotesize},
  best state/.style={state,draw=BookTeal,fill=FigureInputFill,line width=1pt},
  trans/.style={draw=FigureMuted,line width=.5pt,
    -{Stealth[length=1.5mm,width=1mm]}},
  best trans/.style={draw=BookTeal,line width=1.3pt,
    -{Stealth[length=2mm,width=1.4mm]}},
  obs/.style={figure node,minimum height=6mm,minimum width=13mm,
    draw=FigureInput,fill=FigureInputFill},
  value/.style={font=\figurefont\scriptsize,text=FigureInk,fill=white,inner sep=1pt}
]
  % 状态格点：上行为N，下行为V
  \node[best state] (n1) at (0,22) {N};
  \node[state]      (v1) at (0,0)  {V};
  \node[state]      (n2) at (40,22) {N};
  \node[best state] (v2) at (40,0)  {V};
  \node[best state] (n3) at (80,22) {N};
  \node[state]      (v3) at (80,0)  {V};

  % 全部转移（细灰线）
  \foreach \a/\b in {n1/n2,n1/v2,v1/n2,v1/v2,n2/n3,n2/v3,v2/n3,v2/v3}{
    \draw[trans] (\a) -- (\b);
  }
  % Viterbi回溯得到的最优路径
  \draw[best trans] (n1) -- (v2);
  \draw[best trans] (v2) -- (n3);

  % 每个状态旁的δ值与回溯指针
  \node[value,above=1mm of n1] {$\delta_1=0.35$};
  \node[value,below=1mm of v1] {$\delta_1=0.03$};
  \node[value,above=1mm of n2] {$\delta_2=0.014$，$\psi=\mathrm{N}$};
  \node[value,below=1mm of v2] {$\delta_2=0.168$，$\psi=\mathrm{N}$};
  \node[value,above=1mm of n3] {$\delta_3=0.05376$，$\psi=\mathrm{V}$};
  \node[value,below=1mm of v3] {$\delta_3=0.00336$，$\psi=\mathrm{V}$};

  % 观测与时刻
  \node[obs] (o1) at (0,-17) {会议};
  \node[obs] (o2) at (40,-17) {改至};
  \node[obs] (o3) at (80,-17) {周五};
  \node[figure note,below=1mm of o1] {$t=1$};
  \node[figure note,below=1mm of o2] {$t=2$};
  \node[figure note,below=1mm of o3] {$t=3$};

  % 图例
  \draw[best trans] (100,16) -- (110,16);
  \node[anchor=west,font=\figurefont\scriptsize] at (111,16) {最优路径};
  \draw[trans] (100,8) -- (110,8);
  \node[anchor=west,font=\figurefont\scriptsize] at (111,8) {其他转移};
\end{tikzpicture}
